\section{Local causal construction and derived positive service}
\label{sec:local}
\subsection{The moving cutoff does not require a positive old rate}
Compare two histories $(g_i,u_i,Y_i)$ with identical past through $t_0$,
$0\leq g_i\leq G:=R/c_0$ and bounded marks. For fixed $t>t_0$, put
\begin{equation}
 \epsilon=\sup_{s\leq t}|Y_1(s)-Y_2(s)|
 \leq\int_{t_0}^t|g_1-g_2|\dd s,\qquad
 |a_{m,1}(t)-a_{m,2}(t)|\leq\epsilon. \label{eq:cutoffdiff}
\end{equation}
In \eqref{eq:renewaltime}, first hold one indicator fixed. Bounded $D$
and its mark Lipschitz constant give
$R\int|\Delta g|+k_jG\int|\Delta u|$. Where the indicators differ,
$|Y_2(s)-a_{m,2}(t)|\leq2\epsilon$. The $g_2\,ds$ mass in this rank
band is at most $4\epsilon$, including any zero-rate plateaus. Hence
\begin{equation}
 |\Delta\loss(t)|\leq\frac1m\left(
 5R\int_{t_0}^t|\Delta g_E|\dd s+
 k_jG\int_{t_0}^t|\Delta u_E|\dd s\right). \label{eq:volterra}
\end{equation}
The argument also covers recorded old mark jumps. It avoids differentiating
a rank inverse or assuming an unphysical lower bound on E flow.

On a valid selected node regime, \eqref{eq:rootderiv} and bounded $U'$
on $[0,R]$ give local uniform Lipschitz constants $L_g,L_u$ for the E
flow and its mark as functions of $\budget$. The budget map is 1-Lipschitz.
On a sufficiently short interval with the same past,
\begin{equation}
 \|\Delta\mathcal T\budget\|_\infty
 \leq\frac{\tau}{m}(5RL_g+k_jGL_u)\|\Delta\budget\|_\infty.
 \label{eq:contraction}
\end{equation}
Choose the proof interval short enough for contraction and to remain in
a valid chart. Bounded increments give causal construction and continuous
dependence until its next guard. This is a proof device, not a simulator
time step. It does not by itself establish the validity or uniqueness of
every queue switch; Section~\ref{sec:global} joins the selected regimes.

\subsection{Positive budget is a consequence, not an added floor}
Admitted E rate obeys $g_E\leq G$. Since $D(\theta;u)\leq D(\theta;0)$
decreases with age, the maximal possible account fills the youngest ages
at rate $G$ up to mass $m$. Rearrangement gives
\begin{equation}
 \loss\leq p_{\max}:=\frac Gm\int_0^{m/G}D(\theta;0)\dd\theta<R,
 \qquad\budget\geq B_{\min}:=R-p_{\max}>0. \label{eq:budgetfloor}
\end{equation}
Strictness follows from $a_r>0$, making $D(\theta;0)<R$ for positive
age. The bound also holds before the window fills. It is common across
advice at fixed parameters, but can be very small. Clipping in
\eqref{eq:renewal} is inactive in this domain.

From \eqref{eq:renewalupdate}, the retained mass is at most $m$, each
boundary term is bounded by $GR/m$ and $|D_\theta|\leq a_rk_j$.
Therefore the account is Lipschitz with
\begin{equation}
 |\dot\budget|\leq L_B:=2GR/m+a_rk_j\quad\hbox{a.e.}
 \label{eq:budgetrate}
\end{equation}
This allows mark jumps at queue events. Also the renewal integrand is a
subset of the nonnegative all-history additive integrand, so Tonelli gives
\begin{equation}
 \int_0^\infty\loss(t)\dd t\leq
 \overline M_E W_0/m,\quad
 \overline M_E=P_E(L/v)\leq M,\quad W_0=W(0).
 \label{eq:integratedloss}
\end{equation}
It is a bound, not adoption of the additive rival.

\subsection{Two retained limits}
For a constant free A offer $a<R$, $g_E\leq(\budget-a)/c_0$ and
\eqref{eq:renewalupdate} gives $\dot\loss\leq Rg_E/m$. With
$x=\budget-a$, $\dot x\geq-Rx/(mc_0)$, hence
\begin{equation}
 x(t)\geq(R-a)e^{-Rt/(mc_0)}>0. \label{eq:constantA}
\end{equation}
The selected smooth E source cannot generate a first A queue in finite
time in that constant-A limit. H1's A arrivals increase, adding
$-\dot a$ to $\dot x$; Section~\ref{sec:witness} constructs A delay there.

For a constant E rate and mark across a filled renewal window,
\begin{equation}
 \loss=\frac{g_E}{m}W_{m/g_E}(u_E),\quad
 W_h(u)=\int_0^hD(\theta;u)\dd\theta,\quad
 g_A+c_Eg_E=R-\frac{g_E}{m}W_{m/g_E}(u_E). \label{eq:stationaryrenewal}
\end{equation}
This is a restricted stationary target for a simpler instantaneous rival,
not an exact transient reduction or proof that history is necessary.

\section{Whole-domain events, unique continuation and finite clearance}
\label{sec:global}
The result below is restricted to \eqref{eq:domain}--\eqref{eq:workdomain},
the monotone H1 pulse, neutral initial account and selected causal guards.
It is not a general invariant LTM junction theorem.

\subsection{E has at most one waiting episode in this primary domain}
Initially $\budget=R$ and A's offer is below $R$. Let
\begin{equation}
 W_f(v)=\lambda_A(t_f+)+c_0\lambda_E(t_f+). \label{eq:frontwork}
\end{equation}
If $W_f\leq R$, both free streams pass. Because $R\leq C(v_-)$,
$v\geq u_R$, so free E has zero new loss. The free work
$s_v+\delta\lambda_E$ decreases through the pulse; free passage persists.
If $W_f>R$, E's free service test fails and the queued root satisfies
$g_E\leq(R-\lambda_A)/c_0<\lambda_E$. This is a compatible strict
E birth while A remains free.

At a later first E zero hit during the pulse, if A is free, the pre-hit
queued rate implies
\begin{equation}
 \lambda_E\leq g_E^Q\quad\Longrightarrow\quad
 \lambda_A+c_0\lambda_E\leq\budget. \label{eq:EclearAfree}
\end{equation}
If A is queued, its half-work branch implies
\begin{equation}
 c_0\lambda_E\leq c_0g_E^Q\leq\budget/2. \label{eq:EclearAqueued}
\end{equation}
Thus free E is feasible. Its new profiles are zero; old profiles only age
or retire, making $\budget$ nondecreasing. E's arrival offer keeps
decreasing, so \eqref{eq:EclearAqueued} persists while A queues. At A
clearance its former service $\budget-c_0\lambda_E$ admits joint free
flow, and the joint free workload thereafter decreases. Simultaneous
clearance also gives joint feasibility because both individual half-work
offers fit. After the pulse rear there are no new arrivals. E therefore
cannot re-form after clearing in this domain.

\textbf{Scope is essential.} The proof uses the monotone H1 pulse,
free-mark neutrality and selected recovery/priority law, not finiteness
alone. Later/bursty E arrivals in another setting or native finite cohort
can form another queue; that event is not automatically a failed evaluator.

\subsection{A events cannot accumulate infinitely fast}
While E queues during the pulse, A is born when its free offer first loses
feasibility at $\lambda_A=\budget$. Switching to cooperation gives
$\dot q_A=\budget/2\geq B_{\min}/2>0$. If A later clears before E,
the pre-hit derivative requires $\lambda_A\leq\budget/2$.
For $y=\budget-2\lambda_A$, birth has $y=-\budget\leq-B_{\min}$
and clearance has $y\geq0$. By \eqref{eq:pulse} and
\eqref{eq:budgetrate}, a birth-to-clearance interval is at least
\begin{equation}
 B_{\min}/(L_B+2L_A). \label{eq:birthclear}
\end{equation}
At clearance the free margin $x=\budget-\lambda_A\geq B_{\min}/2$.
A later birth needs $x=0$, taking at least $B_{\min}/[2(L_B+L_A)]$.
A common smaller separation is
\begin{equation}
 \tau_0=\frac{B_{\min}}{2(L_B+2L_A)}>0. \label{eq:eventseparation}
\end{equation}
It is derived, not imposed as a dwell timer. E clearance and the single
pulse rear add finitely many exceptions. Once E is free, a cleared A cannot
re-form because joint free work fits and decreases while service recovers.
After the rear, positive queues only drain. Tangential zero hits use the
same free-admission test and still have the post-clearance margin.

\subsection{Joining the regimes and clearing the cohort}
\begin{theorem}[Conditional complete forward evolution]
For every constant advice in the original interval, under the retained
primary assumptions and causal convention, the selected law has a unique
causal piecewise-regime continuation through full cohort clearance.
\end{theorem}
\begin{proof}
Use the local contraction in Section~\ref{sec:local} on each valid chart.
At the front E birth has the strict compatible sign above. An A birth
has positive right-hand derivative. At A zero hits free admission is
feasible; at E zero hits use \eqref{eq:EclearAfree}--\eqref{eq:EclearAqueued}.
Simultaneous hits admit both free streams. Equality alone retains maximal
feasible free continuation, so an alternative queued root cannot create
an untriggered queue. Flux values at individual switching instants do not
alter cumulative mass.

State and service remain bounded. The finite pulse admits at most one E
episode, and \eqref{eq:eventseparation} prevents infinitely many A switches
on it. After the rear no new queues arise, with at most two clearances.
Thus local continuation joins into one unique causal history. No
transversality condition or reduced advice neighborhood is imposed for
this forward result. Smooth advice dependence is a separate claim.

All mass is eligible by $t_r=L/v$. If $Q=q_A+q_E>0$ thereafter, available
work is used, so
$g_A+g_E\geq\budget/c_{\max}\geq B_{\min}/c_{\max}$.
Consequently
\begin{equation}
 T_{\rm clear}(v)\leq\frac Lv+
 \frac{c_{\max}Q(L/v)}{B_{\min}}
 \leq T_{\max}:=\frac L{v_-}+\frac{c_{\max}M}{B_{\min}}<\infty.
 \label{eq:clearance}
\end{equation}
The far lane clears without a queue. Remaining profiles then age out
within their finite support. This finishes the forward construction.
\end{proof}
Another sufficient clearance bound follows by integrating current work
and \eqref{eq:integratedloss} after the rear:
\begin{equation}
 T_{\rm clear}\leq\frac Lv+
 \frac{c_{\max}Q(L/v)+\overline M_EW_0/m}{R}.
 \label{eq:clearancearea}
\end{equation}
Neither estimate asserts a practical simulation runtime or physical realism
of the sufficient parameter domain.

\section{Conservation and the complete local objective}
\label{sec:accounting}
Combining \eqref{eq:upstreammass} and \eqref{eq:terminalcounts} gives
\begin{equation}
 s(t)\rho_0+\sum_jN_j(t)+\sum_jY_j(t)=M. \label{eq:mass}
\end{equation}
All quantities are continuous traffic mass. Work weights change available
service, never the amount transferred or removed. Retiring a profile
removes information weight, not population.

Full local TST is
\begin{equation}
 J(v)=\int_0^{T_{\max}}\left[M-\sum_jY_j(t)\right]\dd t.
 \label{eq:tst}
\end{equation}
By \eqref{eq:clearance}, using any larger bound adds only zero. Every
initial position $x\in[0,\ell]$ has no-delay eligibility time $(L-x)/v$.
Paired early transfers do not change the total eligibility history. Hence
\begin{align}
 \int_0^\infty\left[M-\sum_jP_j(t)\right]\dd t
 &=\int_0^\ell\rho_0\frac{L-x}{v}\dd x
 =\frac{M(L-\ell/2)}v,\nonumber\\
 M-\sum_jY_j&=M-\sum_jP_j+\sum_jq_j,\nonumber\\
 J(v)&=\frac{M(L-\ell/2)}v+\int_0^\infty(q_A+q_E)\dd t.
 \label{eq:tstqueue}
\end{align}
The far lane has $q_B=0$. Both A and E waiting are included exactly once.
Eligible waiting area measures excess time; it is not a physical queue
length or a declaration of when a vehicle first reaches a queue tail.

Completion is at the choke. A crossed E unit can retain a causal restriction
for following unfinished traffic, whose waiting is counted. The completed
unit's own post-line travel is excluded. Adding a recovery-zone objective
would be a separate G4 choice; it is not adopted by this accounting.

\section{Initial information, negative controls and sensitivity limits}
\label{sec:information}
Known initial E mass $E_0$ and continuing symmetry reconstruct
\begin{equation}
 e_0=E_0/\ell,\qquad a_0=b_0=(M-E_0)/(2\ell). \label{eq:scalar}
\end{equation}
Together with fixed parameters and empty terminal/account histories,
the forward theorem yields identical predictions for aggregate-plus-$E_0$
and explicit lane state at the same information, for every permitted
advice. Predicted recovery history is not an extra initial measurement;
the current meter reading alone is not sufficient dynamic state.

The archived predecision aggregate-history collision remains upstream of
an empty terminal and is unaffected by the terminal substitution. Its
reachability/observation assumptions remain conditional. No action split
follows solely from that ambiguity: the amended objectives still require
their own comparison. Strong aggregate/history, scalar, model-form and
bounded-search rivals remain in scope.

If $\delta=0$, the front free work is $s_v<R$. Equal free marks make the
mismatch term zero; no queue or recovery activates. Therefore
\begin{equation}
 J(v)=M(L-\ell/2)/v \label{eq:negativecontrol}
\end{equation}
and fastest advice is common for every preparation. No E is another
neutral limit. This is a property of the selected ordinary-underload
member, not a universal explanation of congestion.

With $z_f=L_p-\ell$ and $e_f(v)=e(z_f;v)$, H1 comparison makes $e_f$
nondecreasing in $v$. Thus
\begin{equation}
 W_f(v)=v\{r+\delta e_f(v)\} \label{eq:frontthreshold}
\end{equation}
is strictly increasing for $r>0$. There is at most one initial queue
threshold; it can lie outside the advice interval. Equality retains free
passage. At an interior threshold with $a=\lambda_A>0$,
$e=\lambda_E>0$, $a+c_0e=R$ and $\chi>0$,
\begin{equation}
 \He(e;a,v)-(R-a)=\chi a e\{v-U(e)\}>0. \label{eq:fluxgap}
\end{equation}
The queued-root limit is below the free E flow, so local flux/marks may
change abruptly on activation. This is not a proof that full $J$ jumps.
Forward uniqueness does not establish smooth derivatives across it.

Within regular charts, \eqref{eq:localbenefit} and ordinary implicit
derivatives describe local response. A full-objective derivative must
also include H1 branch changes, action-dependent arrival times, recovery
marks/aging/retirement and queue events through complete clearance. This
report does not execute or finish that remaining gradient review; it
records why S13's simpler derivative cannot be reused. No AD, numerical
finite differences, policy optimization or useful-gradient claim follows.
